\begin{statementbox}
设 $f(x)=Ax+\frac Bx$，其中 $A,B>0$，$x\ne0$。
\begin{center}
\begin{minipage}{0.98\linewidth}
\centering
\textbf{两支曲线，留下中间的空档}\par\medskip
\begin{tikzpicture}[font=\small,>=Stealth]

\begin{scope}[x=1.08cm,y=0.63cm]
\fill[gray!10] (-3.8,-2) rectangle (3.8,2);
\draw[->] (-3.9,0)--(4,0) node[right] {$x$};
\draw[->] (0,-4.7)--(0,4.8) node[above] {$y$};
\draw[CStatement,very thick,domain=0.24:3.6,samples=100] plot (\x,{\x+1/\x});
\draw[CStatement,very thick,domain=-3.6:-0.24,samples=100] plot (\x,{\x+1/\x});
\draw[dashed] (-3.6,2)--(3.6,2) (-3.6,-2)--(3.6,-2);
\fill[CStatement] (1,2) circle (2.4pt) (-1,-2) circle (2.4pt);
\node[above right,fill=white,inner sep=2pt] at (1,2) {$(1,2)$ 最低点};
\node[below left,fill=white,inner sep=2pt] at (-1,-2) {$(-1,-2)$ 最高点};
\node[fill=white] at (2.1,0.8) {$-2<y<2$ 取不到};
\node[above left] at (-1.1,4.1) {$y=x+\frac1x$};
\node[below left] at (0,0) {$O$};
\draw[very thick,CStatement,->] (4.8,2)--(4.8,4.6);
\draw[very thick,CStatement,->] (4.8,-2)--(4.8,-4.6);
\fill[CStatement] (4.8,2) circle (2.4pt) (4.8,-2) circle (2.4pt);
\node[right] at (4.8,2) {$2$}; \node[right] at (4.8,-2) {$-2$};
\node[above] at (4.8,4.8) {值域};
\end{scope}
\end{tikzpicture}\par\smallskip
图取 $A=B=1$；一般情形的两端为 $\pm2\sqrt{AB}$，横坐标为 $\pm\sqrt{B/A}$。实心点表示端点能取到。
\end{minipage}
\end{center}


	\textbf{\textcolor{CStatement}{条件}}
	\begin{description}[style=nextline, leftmargin=2.8em, labelsep=0.5em, itemsep=0.35em]
		\item[\textbf{\textcolor{CStatement}{条件1（参数及定义域：）}}] $A>0$, $B>0$, $x\neq0$.
	\end{description}

	\textbf{\textcolor{CStatement}{结论}}
	\begin{description}[style=nextline, leftmargin=2.8em, labelsep=0.5em, itemsep=0.45em]
		\item[\textbf{\textcolor{CStatement}{结论一（正数最小值）：}}]
			\textbf{\textcolor{CStatement}{直观描述：}}
			当 $x>0$ 时，函数有最小值 $2\sqrt{AB}$.
			\[
				f(x)\ge 2\sqrt{AB},\quad x>0,
			\]
			等号成立当且仅当 $x=\sqrt{\frac{B}{A}}$.

		\item[\textbf{\textcolor{CStatement}{结论二（负数最大值）：}}]
			\textbf{\textcolor{CStatement}{直观描述：}}
			当 $x<0$ 时，函数有最大值 $-2\sqrt{AB}$.
			\[
				f(x)\le -2\sqrt{AB},\quad x<0,
			\]
			等号成立当且仅当 $x=-\sqrt{\frac{B}{A}}$.

		\item[\textbf{\textcolor{CStatement}{结论三（值域并集）：}}]
			\textbf{\textcolor{CStatement}{直观描述：}}
			函数在 $\mathbb{R}\setminus\{0\}$ 上的取值构成两段区间.
			\[
				(-\infty,-2\sqrt{AB}\,]\cup[\,2\sqrt{AB},+\infty).
			\]

		\item[\textbf{\textcolor{CStatement}{推广结论（分式化对勾）：}}]
			\textbf{\textcolor{CStatement}{直观描述：}}
			当 $g\ne0$ 时，通过线性换元可将形如 $\frac{dx^2+ex+f}{gx+h}$ 的分式函数化为对勾函数形式再求值域.
			\[
				\frac{dx^2+ex+f}{gx+h}\;\xrightarrow{t=gx+h}\; At+\frac{B}{t}+C,
			\]
			如果换元后的系数 $A,B$ 仍满足 $A>0,B>0$，并且新变量 $t$ 的范围允许相应取等点，则可对 $At+\frac{B}{t}$ 部分应用上述最值结论，再加常数 $C$；若系数正负或 $t$ 的范围不符合条件，不能直接套公式，应另外讨论.
	\end{description}

	\textbf{\textcolor{CStatement}{等号/取等条件：}}
	在 $x>0$ 时，基本不等式 $Ax+\frac{B}{x}\ge 2\sqrt{AB}$ 等号成立当且仅当 $Ax=\frac{B}{x}$，即 $x=\sqrt{\frac{B}{A}}$；在 $x<0$ 时，$Ax+\frac{B}{x}\le -2\sqrt{AB}$ 等号成立条件同样是 $Ax=\frac{B}{x}$，解出 $x=-\sqrt{\frac{B}{A}}$. 两种情形均由均值不等式取等条件决定.


\medskip
\textbf{\textcolor{CStatement}{结论四（几何应用：反比例函数两支的最短距离）}}
\begin{center}
\begin{minipage}{0.98\linewidth}
\centering
\textbf{两支被两条平行线隔开：最近连线垂直穿过}\par\medskip
\begin{tikzpicture}[font=\small,>=Stealth]

\begin{scope}[x=1.05cm,y=1.05cm]
\fill[CProof!6] (-3,1)--(-1,3)--(3,-1)--(1,-3)--cycle;
\draw[->,gray] (-3.3,0)--(3.4,0) node[right] {$x$};
\draw[->,gray] (0,-3.3)--(0,3.4) node[above] {$y$};
\draw[CStatement,very thick,domain=0.32:3.1,samples=100] plot (\x,{1/\x});
\draw[CStatement,very thick,domain=-3.1:-0.32,samples=100] plot (\x,{1/\x});
\draw[CProof,thick] (-1,3)--(3,-1);
\draw[CProof,thick] (-3,1)--(1,-3);
\draw[CExample,very thick] (-1,-1)--(1,1);
\fill[CExample] (1,1) circle (2.3pt) (-1,-1) circle (2.3pt);
\node[right,fill=white] at (1.1,1.1) {$(1,1)$};
\node[left,fill=white] at (-1.1,-1.1) {$(-1,-1)$};
\node[text=CProof] at (2.1,2.1) {$x+y=2$};
\node[text=CProof] at (-2.1,-2.1) {$x+y=-2$};
\draw (0.86,0.86)--(0.72,1)--(0.86,1.14);
\end{scope}
\end{tikzpicture}\par\smallskip
图取 $k=1$，强调线段长 $2\sqrt2$。一般情形：两条线为 $x+y=\pm2\sqrt k$，间距 $2\sqrt{2k}$，最近点为 $(\pm\sqrt k,\pm\sqrt k)$。
\end{minipage}
\end{center}

设 $k>0$，在反比例函数 $y=\frac{k}{x}$ 的第一象限分支取点 $P$，第三象限分支取点 $Q$，则
\[
\boxed{PQ\ge 2\sqrt{2k}},\qquad
\boxed{PQ_{\min}=2\sqrt{2k}}.
\]
等号成立当且仅当
\[
P(\sqrt{k},\sqrt{k}),\qquad Q(-\sqrt{k},-\sqrt{k}).
\]
特别地，当 $k=1$ 时，$PQ_{\min}=2\sqrt2$，两点为 $(1,1)$ 与 $(-1,-1)$.
\par\textbf{与本节主结论的联系：} 反比例函数 $y=\frac{k}{x}$ 本身\textbf{不是}标准的 $Ax+\frac Bx$（$A,B>0$）形式；但曲线上一点的坐标和 $x+y=x+\frac{k}{x}$ 正是对勾函数，其最小值 $2\sqrt{k}$ 可以用来估计两支间的距离。这里 $k>0$ 是两支分别落在第一、第三象限的必要条件。

\end{statementbox}
